\emph{Step 1 (the generator).} By Proposition \claimref{field-theory.where-each-lives}, $I_{6} = I_{4}E$ with $E^{2} = 1$ and $[I_{4}, E] = 0$; $E$ is the product of two vectors orthogonal to spacetime, so it commutes with $\sigma$ and with $e_{2}e_{3}$. Hence $I_{6}\sigma = I_{4}E\sigma = E\,I_{4}\sigma$, and $I_{4}\sigma = e_{0}e_{1}e_{2}e_{3}e_{0}e_{1} = (e_{0}e_{1})^{2}e_{2}e_{3} = e_{2}e_{3}$ (Proposition \claimref{state-selection.euclidean-partner-curvature-selected}). So $I_{6}\sigma = E\,e_{2}e_{3}$, the blade $e_{+}e_{-}e_{2}e_{3}$, and $(E\,e_{2}e_{3})^{2} = E^{2}(e_{2}e_{3})^{2} = -1$. It commutes with $\sigma$ because $I_{6}$ is central in the even algebra and $\sigma$ is even. The exponential of an element of square $-1$ is $\cos(\theta/2) + \sin(\theta/2)\,I_{6}\sigma$, which is $-1$ at $\theta = 2\pi$ and $+1$ at $4\pi$, as in Proposition \claimref{singularity.euclidean-rotor-where-it-lies}(1). By Fact \claimref{odd-extensions.imported-conformal-completion} $I_{6}$ is the complex unit of $\CS$. The store's conventions define that unit as right multiplication by $I_{6}$; on $\CS$ left and right multiplication by $I_{6}$ agree, because $I_{6}$ is central in $\Cl^{+}(2,4)$. So on $\CS$ this element is the continued rotor $R_{I_{6}\theta}$, the unit of the continuation read as $I_{6}$.

\emph{Step 2 (the splitting).} $\varepsilon_{\pm} = \tfrac12(1 \pm E)$ are orthogonal idempotents summing to $1$ and commute with $e_{2}e_{3}$, and $\varepsilon_{\pm}E = \pm\varepsilon_{\pm}$. So $\varepsilon_{\pm}R^{W}_{\theta} = \varepsilon_{\pm}(\cos(\theta/2) \pm \sin(\theta/2)e_{2}e_{3}) = \varepsilon_{\pm}R_{\pm\theta}$, and summing gives the displayed identity. The identity holds in the algebra, hence on every left module. On $\CS$ the eigenspaces $L_{0}$, $L_{\infty}$ of $E$ are $\varepsilon_{+}\CS$ and $\varepsilon_{-}\CS$, and on them $I_{6} = \pm I_{4}$ (Proposition \claimref{field-theory.where-each-lives}(3)), which gives (ii).

\emph{Step 3 (not a rotor).} The reverse of a $4$-blade is itself, so $\rev{R^{W}_{\theta}} = R^{W}_{\theta}$ and $R^{W}_{\theta}\rev{R^{W}_{\theta}} = (R^{W}_{\theta})^{2} = R^{W}_{2\theta} = \cos\theta + \sin\theta\,E\,e_{2}e_{3}$, which is $1$ only for $\theta \in 2\pi\Z$. For a vector $x$, $(E e_{2}e_{3})x(E e_{2}e_{3})^{-1}$ is a vector, but the cross term $\cos(\theta/2)\sin(\theta/2)\,[E e_{2}e_{3}, x]$ of the conjugate of $x = e_{2}$ is $2\cos(\theta/2)\sin(\theta/2)\,E\,e_{3}$ up to sign, of grade three, so the conjugate is not a vector unless $\cos(\theta/2)\sin(\theta/2) = 0$. An isometry of the ambient, or a holonomy of a real spin connection, is implemented by an element of the Clifford group, so $R^{W}_{\theta}$ is neither; the angular rotor $R_{\theta}$ is a rotor and fixes $e_{0}$ and $e_{1}$. This is (iii).

The script \file{check.py} verifies in exact arithmetic in $\Cl(2,4)$ the identities of Step 1, the idempotents, the splitting of (ii) at four rational points of the circle, the inverse, the failure of $R\rev{R} = 1$ and of vector preservation at those points, the rotor property of $R_{\theta}$ and its fixing of $e_{0}$, $e_{1}$, and the periods at $\theta = \pi, 2\pi, 4\pi$.
